\begin{claim}{C.43}{Proved}{The lifted boundary quotient is a product projection}
For the actual quotient and covering maps in C.42, every
fiber of the real lift
$F:\widetilde O\to\R$ is connected.
There is a homeomorphism
\[
 T:\R\times\R\longrightarrow\widetilde O,
 \qquad F(T(u,t))=t.                                    \tag{C.43.1}
\]
In particular the components in C.42 form a foliation
with genuine local product charts. The original map
$\Phi_z:O\to B$ also has local product charts with its
specified target coordinate. The construction treats
one quotient map at a time; its auxiliary flow and
the homeomorphism $T$ carry no asserted equivariance
under the original operator family.
\end{claim}

\begin{proof}
Write $X=\widetilde O$ and let $\mathcal L$ be the
closed Polish group (C.42.4), acting transitively on
the covering plane $X$. For $v\in X$, let $C_v$ be
the connected component of $F^{-1}(F(v))$ containing
$v$. By C.42 these sets partition $X$ into properly
embedded lines. The equation $Fa=bF$ for an element
$(l,a,b)\in\mathcal L$ implies
\[
                         a(C_v)=C_{a(v)}.                \tag{C.43.2}
\]
Thus this is the actual component partition preserved
by the given compatible action.

\emph{Regularity and local products.}
Let $I\subset C_v$ be a compact arc containing $v$,
and let $\epsilon>0$. In a fixed metric on the
covering plane, the condition
\[
       \sup_{w\in I}d(a(w),w)<\epsilon
\]
defines an identity neighborhood in $\mathcal L$.
The Polish open orbit-map theorem makes its image
at $v$ a neighborhood of $v$. Consequently, for
every sufficiently near $v'$, some such $a$ sends
$v$ to $v'$. Restricting $a$ to $I$ gives an
$\epsilon$-close embedding of $I$ into $C_{v'}$
carrying $v$ to $v'$. This verifies precisely the
regularity condition of Aarts--Oversteegen
\cite[Definition 5.3]{C_AartsOversteegen1991}.
Their local-section and flowbox theorems
\cite[Theorems 5.6--5.7]{C_AartsOversteegen1991}
therefore provide open product neighborhoods
$S\times(-1,1)\to X$, with each vertical plaque
contained in one $C_v$.

The transverse factor is an interval locally.
Indeed it is metrizable, and its product with
an interval is a surface without boundary.
The elementary product-cancellation argument
proved in C.33 applies: local path connectedness
descends to $S$, embedded open arcs in $S$ are
open by invariance of domain, and an endpoint
would have zero local second homology in the
product. Thus $S$ is locally a one-manifold
without boundary. Shrink to an interval and
write one of these charts as
\[
 \chi:J\times(-1,1)\longrightarrow V\subset X.
\]
Because every plaque lies in a fiber,
$F\chi(s,u)=f(s)$ for a continuous function $f$.
For each open $W\subset J$, the set
$\chi(W\times(-1,1))$ is open in $X$, so the
openness of $F$ proves that $f(W)$ is open.
A continuous open map between intervals is
strictly monotone: failure of strict monotonicity
produces a local maximum, local minimum, or a
constant interval, each contradicting openness.
It follows that $f$ is a homeomorphism onto
an open interval. Replacing the transverse
coordinate by $f(s)$ gives actual charts
in which $F$ is the projection.

\emph{The component quotient and its topology.}
Let $Y$ be the set of components $C_v$, equipped
with the quotient topology, and let
$\pi:X\to Y$ be the quotient map.
This map is open. To verify this directly,
let $U\subset X$ be open and let $v$ belong
to the saturation of $U$. Choose
$w\in C_v\cap U$ and the compact arc in $C_v$
joining $v$ to $w$. Apply the regularity
argument above with an error small enough
that the image of $w$ remains in $U$.
For every $v'$ in a neighborhood of $v$
the resulting image arc joins $v'$ to a
point of $U$ in $C_{v'}$. The saturation
of $U$ is therefore open, which proves
the assertion about $\pi$.

In a chart where $F$ is projection, distinct
plaques have distinct $F$-values. They therefore
belong to distinct global components.
The transversal maps homeomorphically onto
$\pi(V)$: its inverse is the restriction
of the continuous factor map $\overline F$
in
\[
 X\xrightarrow{\ \pi\ }Y
   \xrightarrow{\ \overline F\ }\R,
 \qquad F=\overline F\pi.                               \tag{C.43.3}
\]
Here $\pi(V)$ is open by the preceding
paragraph, and the forward inverse in
the chosen target coordinate is the
continuous transversal section.
Thus $Y$ is a one-manifold without boundary
in the sense that it is locally an interval;
no Hausdorff property has been assumed.
The same charts show that $\overline F$
is a local homeomorphism.

The quotient $Y$ is connected and second
countable: an open quotient carries a
countable base to a countable base.
Every element of $\mathcal L$ induces a
homeomorphism of $Y$ by (C.43.2), and
transitivity on $X$ makes $Y$ homogeneous.
Baillif--Gabard's theorem
\cite[Theorem 1.1, proved in Section 4]
{C_BaillifGabard2008} now proves that
$Y$ is Hausdorff.
The topological hypotheses here are exact:
second countability gives Lindel\"ofness,
and the actual quotient charts give the
local Euclidean structure required in
that theorem.

One can also see the separation mechanism
directly from their proof. In a locally
Hausdorff space, Zorn's lemma supplies
a maximal Hausdorff open set through
each point; maximality makes it dense.
A countable subcover exists by
Lindel\"ofness. A manifold is locally
Baire and hence Baire, so those dense
open sets have a common point. That
point is separated from every other
point inside one member of the subcover.
Homogeneity transports this property
to every point.

A connected second-countable Hausdorff
one-manifold without boundary is a line
or a circle. The circle alternative is
impossible because $\overline F$ is open
and has image $\R$ by C.42. Thus $Y$
is a line. Its local homeomorphism
$\overline F:Y\to\R$ is globally
strictly monotone: the local direction
cannot change on a connected line.
It is onto, hence a homeomorphism.
Each fiber of $F$ is consequently one
of the proper lines $C_v$.

\emph{An exact global product for this map.}
Orient each leaf using the orientation
of the covering plane together with
the increasing transverse coordinate
$F$. On an overlap of two
orientation-preserving charts $(u,F)$,
the leaf-coordinate transition is
increasing. This supplies the coherent
orientation required for the
Aarts--Oversteegen global-flow theorem
\cite[Main Theorem 0.5, its proof in
Section 4, and Remark 4.4]
{C_AartsOversteegen1991}.
It gives a continuous flow
\[
 R:\R\times X\longrightarrow X
\]
whose orbits are exactly the specified
leaves. The flow is free, since those
leaves are lines rather than circles,
and a nontrivial period would make
the corresponding orbit a circle.

Take any small transversal section
$\sigma:I\to X$ with $F\sigma(t)=t$.
The map
\[
 \Theta_\sigma:\R\times I\longrightarrow F^{-1}(I),
 \qquad (u,t)\longmapsto R(u,\sigma(t))                 \tag{C.43.4}
\]
is continuous and bijective: every
fiber is exactly one flow orbit, and
freeness gives uniqueness of $u$.
Both spaces are surfaces, so invariance
of domain makes (C.43.4) a homeomorphism.
These are actual local trivializations
of a principal $\R$-bundle over $\R$.

For completeness, its global section
can be constructed with the group
coordinate retained. Choose a locally
finite interval cover $(I_i)$ with
sections $\sigma_i$ and a subordinate
partition of unity $(\rho_i)$ whose
supports lie in $I_i$. On overlaps,
let the continuous real transition
functions be defined by
\[
             \sigma_j(t)=R(a_{ij}(t),\sigma_i(t)).
\]
They obey $a_{ij}=a_{ik}+a_{kj}$.
Set $b_i(t)=\sum_j\rho_j(t)a_{ij}(t)$
on $I_i$, extending the individual
terms by zero outside their supports.
The sum is locally finite and
$b_i=a_{ik}+b_k$ on overlaps.
Consequently
$\sigma(t)=R(b_i(t),\sigma_i(t))$
is a well-defined continuous global
section with $F\sigma(t)=t$.
The same map as (C.43.4), now over
all $\R$, gives (C.43.1).

Finally, the identities
$qF=\Phi_zp$ from C.42, together with
local charts for the two covering
maps $p$ and $q$, carry the projection
charts to the original map $\Phi_z$.
This proves its stated local product
property with the actual target
coordinate.
\end{proof}

\paragraph{Scope of the reduction.}
The auxiliary flow in the proof is used
to trivialize one specified quotient.
It has not been shown to commute with
the original operators, preserve another
boundary quotient, or extend across the
frontier of $O$. C.43 therefore strengthens
the structure of the necessary boundary
maps without completing the simultaneous
relative extension in C.12.
